% Section 3.1: finite geometric-Frobenius orbits.
% This file is a chapter fragment; it is \input from
% version/detailed/weil_conjectures_chapter3.tex and must not carry a preamble.

\section{Finite geometric-Frobenius orbits}
\label{sec:weil-finite-orbits}

Throughout this section $\Omega$ is a set and $F\colon\Omega\to\Omega$ is a
self-map; the case of interest is the geometric Frobenius acting on the
$\overline{\Fq}$-points of a variety over $\Fq$, but neither bijectivity of
$F$ nor any geometry is used.  We write $F^{[n]}$ for the $n$-fold iterate of
$F$, so that $F^{[0]}=\mathrm{id}$ and $F^{[m+n]}=F^{[m]}\circ F^{[n]}$.

The goal is \Cref{prop:weil-fixed-points-orbit-count}: the number of points
fixed by $F^{[n]}$ is $\sum_{m\mid n}m\,a_m$, where $a_m$ counts the orbits of
$F$ of exact size $m$; and its Möbius inversion
\Cref{lemma:weil-orbit-count-mobius-numerator-eq}, which recovers $a_m$ from
the fixed-point counts.

\subsection{Periodic points and their orbits}

\begin{definition}[Points of period $n$]
  \label{def:weil-frobenius-fixed-point-set}
  \lean{Function.ptsOfPeriod}
  \leanok

  For $n\in\mathbf N$ put
  \[
    \operatorname{Fix}_F(n):=\{x\in\Omega \mid F^{[n]}(x)=x\},
  \]
  the set of points of period $n$, and write
  $N_n:=\#\operatorname{Fix}_F(n)$ for its cardinality whenever it is finite.
  A point is \emph{periodic} if it lies in $\operatorname{Fix}_F(n)$ for some
  $n\ge 1$; the least such $n$ is the \emph{minimal period} $d_F(x)$.
\end{definition}

\begin{lemma}[Fixed points from the minimal period]
  \label{lemma:weil-frobenius-fixed-iff-orbit-length-divides}
  \uses{def:weil-frobenius-fixed-point-set}
  \lean{Function.isPeriodicPt_iff_minimalPeriod_dvd}
  \leanok

  For every $x\in\Omega$ and every $n\in\mathbf N$,
  \[
    F^{[n]}(x)=x
    \quad\Longleftrightarrow\quad
    d_F(x)\mid n .
  \]
\end{lemma}

\begin{proof}
  \leanok
  This is the defining property of the minimal period: the set of $n$ with
  $F^{[n]}(x)=x$ is closed under subtraction and hence is the set of multiples
  of its least positive element (or all of $\mathbf N$ when $x$ is not
  periodic, in which case $d_F(x)=0$).
\end{proof}

\begin{lemma}[Monotonicity of the fixed-point sets]
  \label{lemma:weil-frobenius-fixed-point-set-monotone}
  \uses{def:weil-frobenius-fixed-point-set,
    lemma:weil-frobenius-fixed-iff-orbit-length-divides}
  \lean{Function.ptsOfPeriod_subset_of_dvd}
  \leanok

  If $m\mid n$ then $\operatorname{Fix}_F(m)\subseteq\operatorname{Fix}_F(n)$.
  In particular $\operatorname{Fix}_F(n)$ finite implies
  $\operatorname{Fix}_F(m)$ finite for every divisor $m$ of $n$.
\end{lemma}

\begin{proof}
  \uses{lemma:weil-frobenius-fixed-iff-orbit-length-divides}
  \leanok
  By \Cref{lemma:weil-frobenius-fixed-iff-orbit-length-divides} membership in
  $\operatorname{Fix}_F(m)$ says $d_F(x)\mid m$, and $d_F(x)\mid m\mid n$.
\end{proof}

\begin{definition}[Orbit of a periodic point]
  \label{def:weil-frobenius-orbit}
  \uses{def:weil-frobenius-fixed-point-set}
  \lean{Function.periodicOrbit}
  \leanok

  For $x\in\Omega$ let
  \[
    \mathcal O(x):=
      \bigl(x,\;F(x),\;F^{[2]}(x),\;\dots,\;F^{[d_F(x)-1]}(x)\bigr)
  \]
  be the orbit of $x$, taken as a cyclic sequence, so that a cyclic shift of
  the listed points is the same orbit.  Its length is $d_F(x)$, and for a
  non-periodic $x$ it is the empty orbit.
\end{definition}

\begin{lemma}[Members of an orbit]
  \label{lemma:weil-frobenius-orbit-members}
  \uses{def:weil-frobenius-orbit}
  \lean{Function.mem_periodicOrbit_iff_exists_lt_minimalPeriod}
  \leanok

  For all $x,y\in\Omega$,
  \[
    y\in\mathcal O(x)
    \quad\Longleftrightarrow\quad
    y=F^{[j]}(x)\ \text{for some}\ 0\le j<d_F(x).
  \]
  No periodicity hypothesis is needed: for a non-periodic $x$ both sides are
  false.
\end{lemma}

\begin{proof}
  \leanok
  Immediate from \Cref{def:weil-frobenius-orbit}, whose underlying list is
  $j\mapsto F^{[j]}(x)$ on $\{0,\dots,d_F(x)-1\}$.
\end{proof}

\begin{lemma}[Orbits are equal or disjoint]
  \label{lemma:weil-frobenius-orbits-equal-or-disjoint}
  \uses{def:weil-frobenius-orbit, lemma:weil-frobenius-orbit-members}
  \lean{Function.periodicOrbit_eq_iff_mem_periodicOrbit}
  \leanok

  Let $x$ be periodic.  Then for every $y\in\Omega$,
  \[
    \mathcal O(y)=\mathcal O(x)
    \quad\Longleftrightarrow\quad
    y\in\mathcal O(x).
  \]
  Consequently two orbits that share a point are equal.
\end{lemma}

\begin{proof}
  \uses{lemma:weil-frobenius-orbit-members}
  \leanok
  If $y\in\mathcal O(x)$ then $y=F^{[j]}(x)$ by
  \Cref{lemma:weil-frobenius-orbit-members}, and iterating $F$ does not change
  the orbit of a periodic point, so $\mathcal O(y)=\mathcal O(x)$.
  Conversely, if $\mathcal O(y)=\mathcal O(x)$ then $y$ is periodic, because
  $\mathcal O(x)$ is nonempty while the orbit of a non-periodic point is
  empty; and a periodic point lies on its own orbit, so
  $y\in\mathcal O(y)=\mathcal O(x)$.
\end{proof}

\begin{lemma}[Periods are constant along an orbit]
  \label{lemma:weil-frobenius-orbit-period-invariant}
  \uses{def:weil-frobenius-fixed-point-set, lemma:weil-frobenius-orbit-members}
  \lean{Function.IsPeriodicPt.of_mem_periodicOrbit}
  \leanok

  If $x\in\operatorname{Fix}_F(n)$ and $y\in\mathcal O(x)$ then
  $y\in\operatorname{Fix}_F(n)$.
\end{lemma}

\begin{proof}
  \uses{lemma:weil-frobenius-orbit-members}
  \leanok
  Write $y=F^{[j]}(x)$ by \Cref{lemma:weil-frobenius-orbit-members}; then
  $F^{[n]}(y)=F^{[j]}(F^{[n]}(x))=F^{[j]}(x)=y$.
\end{proof}

\subsection{Counting orbits}

\begin{definition}[Orbits of a given degree]
  \label{def:weil-degree-frobenius-orbit-set}
  \uses{def:weil-frobenius-orbit}
  \lean{Function.orbitsOfMinimalPeriod}
  \leanok

  For $m\in\mathbf N$ let
  \[
    \mathcal O_m:=
      \{\mathcal O(x) \mid x\ \text{periodic},\ d_F(x)=m\}
  \]
  be the set of orbits of degree $m$, that is, of orbits consisting of exactly
  $m$ points.
\end{definition}

\begin{definition}[Orbit counts]
  \label{def:weil-frobenius-orbit-count}
  \uses{def:weil-degree-frobenius-orbit-set}
  \lean{Function.orbitCount}
  \leanok

  Let $a_m:=\#\mathcal O_m$ be the number of orbits of degree $m$, with the
  convention that an infinite set has count $0$.
\end{definition}

\begin{lemma}[No orbits of degree zero]
  \label{lemma:weil-frobenius-orbit-count-zero}
  \uses{def:weil-degree-frobenius-orbit-set, def:weil-frobenius-orbit-count}
  \lean{Function.orbitsOfMinimalPeriod_zero, Function.orbitCount_zero}
  \leanok

  $\mathcal O_0=\varnothing$ and $a_0=0$.
\end{lemma}

\begin{proof}
  \leanok
  A periodic point has $d_F(x)\ge 1$.
\end{proof}

\begin{proposition}[Fixed points counted orbit by orbit]
  \label{prop:weil-fixed-points-orbit-count}
  \uses{def:weil-frobenius-fixed-point-set,
    lemma:weil-frobenius-fixed-iff-orbit-length-divides,
    lemma:weil-frobenius-orbits-equal-or-disjoint,
    lemma:weil-frobenius-orbit-period-invariant,
    def:weil-frobenius-orbit-count}
  \lean{Function.ncard_ptsOfPeriod_eq_sum_divisors_mul_orbitCount}
  \leanok

  Let $n\ge 1$ and suppose $\operatorname{Fix}_F(n)$ is finite.  Then
  \[
    N_n=\sum_{m\mid n}m\,a_m .
  \]
\end{proposition}

\begin{proof}
  \uses{lemma:weil-frobenius-fixed-iff-orbit-length-divides,
    lemma:weil-frobenius-orbit-members,
    lemma:weil-frobenius-orbits-equal-or-disjoint,
    lemma:weil-frobenius-orbit-period-invariant}
  \leanok
  Partition the finite set $\operatorname{Fix}_F(n)$ into the fibres of
  $x\mapsto\mathcal O(x)$.

  The fibre over $\mathcal O(x)$ is $\mathcal O(x)$ itself: by
  \Cref{lemma:weil-frobenius-orbits-equal-or-disjoint} a point $y$ has
  $\mathcal O(y)=\mathcal O(x)$ exactly when $y\in\mathcal O(x)$, and every
  such $y$ again lies in $\operatorname{Fix}_F(n)$ by
  \Cref{lemma:weil-frobenius-orbit-period-invariant}.  By
  \Cref{lemma:weil-frobenius-orbit-members} that fibre is the image of
  $\{0,\dots,d_F(x)-1\}$ under $j\mapsto F^{[j]}(x)$, and this map is
  injective below the minimal period, so the fibre has exactly $d_F(x)$
  elements.

  Every orbit occurring is of degree dividing $n$, because
  $x\in\operatorname{Fix}_F(n)$ gives $d_F(x)\mid n$ by
  \Cref{lemma:weil-frobenius-fixed-iff-orbit-length-divides}.  Grouping the
  fibres by that degree, the orbits of degree $m$ occurring are precisely all
  of $\mathcal O_m$: conversely, if $d_F(x)=m$ and $m\mid n$ then
  $x\in\operatorname{Fix}_F(n)$.  Summing $d_F(x)=m$ over the $a_m$ orbits of
  degree $m$, and then over the divisors $m$ of $n$, gives
  $N_n=\sum_{m\mid n}m\,a_m$.
\end{proof}

\subsection{Möbius inversion}

\begin{definition}[Möbius numerator]
  \label{def:weil-orbit-count-mobius-numerator}
  \uses{def:weil-frobenius-fixed-point-set}
  \lean{Function.moebiusPtsOfPeriodSum}
  \leanok

  For $m\in\mathbf N$ let
  \[
    M_m:=\sum_{r\mid m}\mu(r)\,N_{m/r}\in\Z ,
  \]
  where $\mu$ is the Möbius function.
\end{definition}

\begin{lemma}[Möbius inversion of the orbit count]
  \label{lemma:weil-orbit-count-mobius-numerator-eq}
  \uses{def:weil-orbit-count-mobius-numerator,
    prop:weil-fixed-points-orbit-count,
    lemma:weil-frobenius-fixed-point-set-monotone}
  \lean{Function.moebiusPtsOfPeriodSum_eq_mul_orbitCount}
  \leanok

  Let $m\ge 1$ and suppose $\operatorname{Fix}_F(m)$ is finite.  Then
  \[
    M_m=m\,a_m .
  \]
\end{lemma}

\begin{proof}
  \uses{prop:weil-fixed-points-orbit-count,
    lemma:weil-frobenius-fixed-point-set-monotone}
  \leanok
  By \Cref{lemma:weil-frobenius-fixed-point-set-monotone} every
  $\operatorname{Fix}_F(k)$ with $k\mid m$ is finite, so
  \Cref{prop:weil-fixed-points-orbit-count} gives
  $\sum_{i\mid k}i\,a_i=N_k$ for every divisor $k$ of $m$.  Möbius inversion
  on the divisor-closed set of divisors of $m$ turns this into
  $\sum_{r\mid m}\mu(r)N_{m/r}=m\,a_m$.
\end{proof}

\begin{corollary}[Integrality]
  \label{cor:weil-orbit-count-mobius-integrality}
  \uses{lemma:weil-orbit-count-mobius-numerator-eq}
  \lean{Function.natCast_dvd_moebiusPtsOfPeriodSum}
  \leanok

  Let $m\ge 1$ and suppose $\operatorname{Fix}_F(m)$ is finite.  Then
  $m\mid M_m$ in $\Z$.
\end{corollary}

\begin{proof}
  \uses{lemma:weil-orbit-count-mobius-numerator-eq}
  \leanok
  Immediate from \Cref{lemma:weil-orbit-count-mobius-numerator-eq}.
\end{proof}

\begin{corollary}[Recovery of the orbit counts]
  \label{cor:weil-orbit-count-mobius-recovery}
  \uses{lemma:weil-orbit-count-mobius-numerator-eq,
    cor:weil-orbit-count-mobius-integrality}
  \lean{Function.moebiusPtsOfPeriodSum_div_natCast}
  \leanok

  Let $m\ge 1$ and suppose $\operatorname{Fix}_F(m)$ is finite.  Then
  \[
    a_m=\frac{1}{m}\sum_{r\mid m}\mu(r)\,N_{m/r}.
  \]
\end{corollary}

\begin{proof}
  \uses{lemma:weil-orbit-count-mobius-numerator-eq}
  \leanok
  Divide \Cref{lemma:weil-orbit-count-mobius-numerator-eq} by $m$, which is
  legitimate in $\Z$ by \Cref{cor:weil-orbit-count-mobius-integrality}.
\end{proof}
